\documentclass[11pt]{article}
\usepackage{biblatex}
\usepackage{amsmath}
\addbibresource{Molecular_Docking_Methods_and_Programs.bib}
\begin{document}
%\chapter*{Molecular Docking Methods and Programs}
\section*{Introduction}
As structure-based drug design developed into a mature science, molecular docking got better with every step. The method depends on how well the conformational search algorithm can find the bioactive conformation(s). However, the search algorithm by itself does not know what a bioactive conformation(s) looks like. By definition, bioactive conformation is the atomic arrangement of a molecule that is suitable for binding to its respective target, and either inhibits/antagonises or activates it. There is always a possibility that one ligand may modulate more than one target. In this case, the bioactive conformation may be different. However, the bioactive conformation is always one of the plausible low-energy states, not necessarily the global minimum. Now that we have established that there are several possible bioactive conformations for a ligand depending upon its binding target, the conformational search algorithm alone is unable to identify it, but it will find it. The scoring function is responsible for scoring and ranking the conformations from the search algorithm. The chapter on scoring functions ((\textbf{see chapter \ref{Scoring_Functions_in_Molecular_Docking}})) will treat this part in depth. Remember that the conformational search tool searches for the bioactive conformation within the receptor binding site; therefore, many low-energy conformations that are observed in pure gas phase or solvent may not be sampled due to clashes with the receptor atoms. This is one of the ways by which docking programs implement a preliminary "bump-check" that can filter out conformations that show clashes with the receptor atoms. This way, more intense calculations are performed on fewer conformations. This improves the speed of docking. 

This chapter will discuss some commonly used docking programs. Few of them need a commercial license for a fee, while are rest are open-source for academic use. 
There are a plethora of docking programs available (\textbf{Table \ref{tab:List}}), some with commercial licenses and others with open-source licenses. Knowing every program to make a choice that best suits you is almost always impossible. However, this is also not necessary as there exists ample literature reporting the use of many of the tools. Irrespective of the choice of the docking program, the general steps will be same; for example, preparation of the receptor (\textbf{see chapter \ref{receptors}}), defining the binding site as search space(\textbf{see chapter \ref{binding_sites}}), preparation of ligands (\textbf{see chapter \ref{ligprep}}), specifying the docking parameters, validating the docking procedure, performing the docking calculations, and analysis and dissemination of the results.  In a nutshell, the general principles remain unchanged; however, there are differences in ligand format, docking methods, and scoring algorithms that are well-documented in the user manual. We strongly recommend that, before using any docking algorithm, the user manual must be consulted to understand each parameter's importance. The developers present their report on several validation studies as a part of the software development, which is a great resource to understand the effect of each parameter on the outcome. A literature survey of the program's usage will also be covered in the user manual. 


\begin{table}[ht]
    \centering
    \caption{List of common molecular docking software}
    \label{tab:List}
    \begin{tabular}{|c|c|c|c|}
    \hline
         \textbf{Software} & \textbf{License type} & \textbf{Scoring Function}  \\
    \hline
        AutoDock & Open Source & Force field  \\
        AutoDock-Vina & Open Source & Empirical  \\
        DOCK & Open Source (Academic) & Force field  \\
        GOLD &  Commercial & Multiple options \\
        FlexX &  Commercial & Empirical \\
        Glide &  Commercial & Empirical \\
        FLEXX & Commercial & ?? \\
        SURFLEX & Commercial & ?? \\
    \hline
    \end{tabular}
\end{table}


\section*{Molecular Docking Programs:}
The following sections will describe the working of the commonly available docking program (See Table \ref{tab:List})

\subsection*{DOCK}
The DOCK program is considered the first tool developed for molecular docking in 1982\cite{bib71}. It was developed by the University of California, San Francisco (UCSF) in the group of Irwin Kuntz. It aims to predict the most favourable binding configurations of small molecules (ligands) within the active sites of macromolecular targets, such as proteins. This computational resource simulates molecular interactions to identify low-energy binding poses, making it an essential asset for drug discovery and structure-based design initiatives. Created to tackle rigid-body docking, DOCK utilises a geometric matching algorithm that aligns the ligand with a negative representation of the receptor's binding pocket. Over the years, its functionality has evolved to incorporate features such as force-field-based scoring, real-time optimisation, and algorithms for flexible ligand docking\cite{bib18}. The program employs scoring functions to assess and rank potential binding poses according to their predicted binding affinities. It accommodates various input file formats, including PDB, MOL2, and SDF, and is extensively utilised for virtual screening and lead optimisation in drug development. DOCK continues to be a fundamental tool in computational chemistry, renowned for its efficiency in exploring molecular interactions and its ability to adapt to new advancements in docking techniques.

\subsubsection*{Methodology:} The docking method described in the sources involves a geometric approach to explore feasible alignments of ligands and receptors of known structure. We describe the process as shown in the original publication\cite{bib71} as follows:

\begin{enumerate}
    \item \textit{Representation}: The receptor and ligand structures are represented using spheres to fill pockets and grooves on the receptor's surface and the space occupied by the ligand. These spheres are generated based on the molecular surface, which is defined using a probe sphere to determine contact and re-entrant points. The spheres are constructed such that each sphere touches the molecular surface at two points and has its center on the surface normal. The number of spheres is reduced by keeping only the smallest radius spheres, filtering based on the angle formed by the two surface points and the sphere centre, and by only retaining one sphere per atom. Receptor spheres are grouped into potential binding sites based on whether they overlap.
   
    \item \textit{Matching}: The ligand spheres are paired with the receptor spheres by comparing internal distances. A ligand sphere can be paired with a receptor sphere if the internal distances of all the spheres in the ligand set match all the internal distances within the receptor set, within an error limit. The pairing process starts by systematically pairing each ligand sphere with each receptor sphere and then finding subsequent pairs based on whether their internal distances match. If the number of assignable pairs is less than four, the match is rejected. To reduce computation time, only ligand and receptor spheres whose centres are approximately the same distance from the centroid of each respective representation are accepted.

   \item \textit{Optimisation}: The suggested pairings are explored by rotating the ligand spheres onto the corresponding receptor spheres using a least-squares algorithm. The rotation/translation matrix generated from the spheres is applied to all ligand atoms. The optimisation procedure manipulates the ligand atom coordinates to reduce atom overlaps, ensure hydrogen-bonding partners, and place ligand atoms within receptor spheres. An overlap error is calculated based on the van der Waals' radii of the ligand and receptor atoms. The ligand is moved further into the receptor site by using the receptor sphere centre as a target for each ligand atom. Displacements are calculated for any ligand atom that violates the overlap constraint or is too far from a receptor polar atom, and the ligand is then moved to reduce these errors.
   
\end{enumerate}

\subsubsection*{scoring:}
The output is a set of ligand conformations/poses with scores based on the degree of overlap. Geometrically similar structures are grouped. The method employs a simplified interaction function that accounts for hard-sphere repulsions and hydrogen bonding. It is important to note that this method assumes rigid molecules and requires prior knowledge of the two structures to atomic detail. The calculated structures are intended as starting points for further refinement using molecular mechanics or interactive computer graphics. The approach also has independent applications, such as the location of binding sites and the design of new ligands.

\subsection*{GLIDE}
One of the most popular methods for docking is Glide, which stands for \textbf{G}rid-based \textbf{L}igand \textbf{D}ocking with \textbf{E}nergetics. This tool was developed by Schrödinger Inc. led by Prof. Richard Frienser and David Shaw from Columbia University and D. E. Shaw Research and Development, respectively \cite{glide001}.

Glide employs a systematic, hierarchical approach to ligand conformational filtering and pose selection, designed to efficiently explore the ligand's conformational space and its potential interactions within the receptor's active site while maintaining computational speed. This methodology involves a series of filters that progressively narrow the search space.
\subsubsection*{How does GLIDE search for ligand conformations?}

The detailed description of each step in the GLIDE docking method is given below:
\begin{enumerate}
    \item Receptor Pre-computation and Grid Generation: The protein receptor's shape and properties are precomputed and represented on grids, using different sets of fields for progressively more accurate scoring. A multigrid strategy is used, with finer resolution (down to 0.4 Å) closer to the protein's van der Waals surface, to ensure accuracy where ligand and protein contact. This pre-processing step creates an efficient and accurate representation of the rigid receptor environment, allowing rapid energy and gradient evaluations during subsequent ligand interaction calculations. The varying grid resolution optimises computational cost while maintaining precision in critical binding regions.

    \item Ligand Conformation Generation: Glide first performs an exhaustive conformational search in the ligand's torsion-angle space to generate initial conformations. A heuristic screen rapidly eliminates high-energy conformers by evaluating torsional energy using a truncated OPLS-AA potential and applying an energy cut-off. Ligands are conceptually divided into a rigid "core" region and flexible "rotamer groups" attached by rotatable bonds. Each core conformation is enumerated, and then combined with all possible rotamer states, allowing thousands of effective conformations to be considered. This step aims for a comprehensive sampling of ligand flexibility, ensuring that relevant bioactive conformations are likely to be generated. The core/rotamer model and heuristic screening make the exhaustive search computationally feasible for even large and flexible ligands. 

     \item Initial Screening of Ligand Poses: For each generated ligand conformation, an exhaustive search of possible positions and orientations is performed within the active site. This multi-stage process includes: 
     \begin{enumerate}
     \item Site Point Selection: Placing the ligand centre at "site points" on a 2{\AA} grid if there's a good match between ligand and receptor surface distances.
     \item Orientation Screening: Testing pre-specified orientations of the "ligand diameter" and rejecting those with excessive steric clashes.
     \item Subset Scoring: Scoring only atoms capable of forming hydrogen bonds or metal interactions (a "subset test").
     \item Full Interaction Scoring: Scoring all ligand-receptor interactions for promising orientations.
     \item These tests use a discretised ChemScore version, rapidly identifying favourable hydrophobic, hydrogen-bonding, and metal-ligation interactions while penalising clashes. Approximately 5000 top greedy-scoring poses are then subjected to a rigid refinement allowing 1{\AA} Cartesian movement.
    \end{enumerate}
    This is a rapid pre-screening stage that drastically reduces the conformational and positional search space. By using a fast, approximate scoring function ("greedy scoring"), it efficiently filters out unlikely poses, making more computationally intensive subsequent steps feasible. The exhaustive nature aims to avoid missing key regions of phase space.
    
    \item Energy Minimization (SP Minimization): A select number of the best poses (typically 400) from the initial screening are subjected to energy minimization using precomputed OPLS-AA van der Waals and electrostatic grids. The minimization often begins on "smoothed" grids to avoid deep local minima, gradually transitioning to the full OPLS-AA non-bonded energy surface ("annealing"). This step accounts for rigid-body translations and rotations, and for flexible docking, also includes torsional motion around rotatable bonds. This step refines the rough poses into more energetically stable configurations within the receptor's force field. The annealing approach helps to overcome local energy barriers, leading to more accurate pose prediction.
    
    \item Post-Docking Minimisation and Monte Carlo Refinement: The very best poses (typically 3 to 6) from the energy minimisation are subjected to an additional Monte Carlo procedure. This procedure explores nearby torsional minima, refining the orientation of peripheral groups and occasionally adjusting internal torsion angles. This final refinement step ensures that the most promising poses are thoroughly optimized, capturing subtle conformational adjustments that may be crucial for accurate binding mode prediction.

    \item Pose Selection (Emodel): The ultimate selection of the best-docked pose is made using Emodel, a composite energy score. Emodel combines the GlideScore (an expanded ChemScore-like empirical function), the ligand-receptor molecular mechanics interaction energy, and the ligand's internal strain energy (for flexible docking). Emodel provides a robust and accurate criterion for selecting the most representative binding pose. By integrating multiple energy components, it offers a more balanced assessment of pose quality compared to using individual scoring functions alone.

\end{enumerate}

\subsubsection*{How does GLIDE Score the docked conformations?}
Glide employs several scoring functions, primarily GlideScore (with Standard Precision (SP) and Extra Precision (XP) versions) and Emodel, a composite score used for final pose selection. These functions integrate various physical and empirical terms to evaluate ligand binding affinity and rank poses.
The general form of the GlideScore 2.5 SP function is given as:
\begin{equation} \label{eq:GLIDEscore}
\begin{split}
\ \Delta G_{\text{bind}} = \\ C_{\text{lipo-lipo}}{\sum}f(r_{lr}) + \\ 
                    C_{hbond-neut-neut}{\sum}g({\Delta}r) h({\Delta}R) + \\
                    C_{hbond-neut-charged}{\sum}g({\Delta}r) h({\Delta}R) + \\ 
                    C_{hbond-charged-charged}{\sum}g({\Delta}r) h({\Delta}R) + \\
                    C_{max-metal-ion}{\sum}f(r_{lm}) +  \\
                    C_{rotb}H_{rotb} + C_{polar-phob}V_{polar-phob} +  \\
                    C_{coul}E_{coul} + C_{vdW}E_{vdW} + solvation terms 
\end{split}
\end{equation}
Here are the key components of Glide's scoring functions:
\begin{enumerate}
    \item GlideScore (General Components - Shared between SP and XP)
    \begin{enumerate}
     \item Lipophilic-Lipophilic Term: This term, similar to that in ChemScore, accounts for interactions between lipophilic atoms of the ligand and the receptor. It heuristically represents the displacement of water molecules from hydrophobic regions by lipophilic ligand atoms, which lowers the free energy via the hydrophobic effect.
     \item Hydrogen-Bonding Term: Based on ChemScore, this term is separated into components with different weights depending on the charge state of the donor and acceptor atoms: neutral-neutral, neutral-charged, and charged-charged. In the optimized scoring function, neutral-neutral contributions are found to be the most stabilizing, while charged-charged are the least important. The quality of hydrogen bonds is evaluated based on geometric criteria (angles and distances).
     \item Metal-Ligand Interaction Term: This term considers interactions with anionic acceptor atoms, recognising a strong preference for coordination of anionic ligand functionality to metal centres in metalloproteases. Glide 2.5 counts only the single best interaction if multiple metal ligations are found. The preference for anionic ligands is suppressed if the metal ion in the apo-protein is neutral.
     \item Coulomb and van der Waals (vdW) Interaction Energies: These molecular mechanics components contribute to GlideScore 2.5. To make the gas-phase Coulomb interaction energy a better predictor, Glide reduces the net ionic charge on formally charged groups (e.g., carboxylates, guanidiniums) by approximately 50\%. And reduces the vdW interaction energies for the directly involved atoms. The full interaction energy is used only for anionic ligand-metal interactions.
      \item Solvation Terms / Desolvation Penalties: Glide 2.5 incorporates a solvation model to address the requirement that charged and polar groups of both the ligand and protein must be adequately solvated. Explicit water molecules are docked into high-scoring poses, and empirical scoring terms measure the exposure of various groups to these waters. Penalties are assessed if a polar or charged ligand or protein group is inadequately solvated or if water molecules are trapped in hydrophobic pockets. This "water-scoring" technology is made efficient by grid-based algorithms.
    \end{enumerate}
    \item  Extra Precision (XP) GlideScore Specific Novel Terms: The XP scoring function and docking protocol were developed to reproduce experimental binding affinities and yield quality enrichments, highlighting the importance of novel molecular recognition and water scoring in separating active and inactive ligands and avoiding false positives. These terms aim to identify specific structural motifs that provide large contributions to enhanced binding affinity.
     \begin{enumerate}
    \item Hydrophobic Enclosure (Ehyd•enclosure): This term rewards situations where groups of lipophilic ligand atoms are enclosed on opposite faces by lipophilic protein atoms. It goes beyond a simple atom-atom pair term, recognizing that enclosed hydrophobic micro-environments are particularly unfavourable for water molecules, leading to significant binding free energy gains upon water displacement. A minimum group size of three connected lipophilic ligand atoms is considered for this effect. The score is capped at 4.5 kcal/mol.
    \item Special Neutral-Neutral Hydrogen-Bond Motifs (Ehb•nn•motif): These motifs receive additional binding affinity increments. They are identified in situations where a water molecule forming a hydrogen bond to the protein would have particular difficulty in making its complement of additional hydrogen bonds, especially when there are hydrophobic protein atoms on two faces. This is often crucial for high potency in pharmaceutical targets like kinases. The reward is 1.5 kcal/mol for a single bond to a ring in hydrophobic environments and 3.0 kcal/mol for a neutral pair in hydrophobic environments. The protein's hydrogen-bond partner must have a limited number of waters in its first or second shell. Correlated hydrogen bonds between protein and ligand can also receive rewards if they are in a challenging hydrophobic environment.
    \item Special Charged-Charged Hydrogen-Bond Motifs (Ehb•cc•motif): These terms signal enhanced binding affinity from charged-charged hydrogen bonds based on:
           \begin{enumerate}
        \item Number of waters surrounding the protein component of the salt bridge, as fully solvent-exposed charged groups are unlikely to participate in enhanced interactions.
        \item Number of charged-charged hydrogen bonds (monodentate, bidentate, or one ligand group to two different protein groups).
        \item Zwitterion ligands: For zwitterions, the long-range effects of solvation free energy are significantly reduced, making the formation of two salt bridges more favourable.
        \item Distinguishing between positive ligand/negative protein and vice versa.
        \item Strength of the electrostatic field at the ligand.
        \item Rewards for various motifs range from 0.5 to 4.7 kcal/mol.
        \end{enumerate}
    \item Pi-Stacking and Pi-Cation Interactions (EPI): These specialised terms, still under development, reward such interactions.
    \item Empirical Correction for Smaller Ligands: An empirical correction enhances the binding affinity of smaller ligands relative to larger ones.
         \end{enumerate}
\item Penalty Terms in XP GlideScore:
    \begin{enumerate}
    \item Desolvation Penalties (E$_{desolv}$): As described under solvation terms above, these specifically penalize inadequately solvated polar or charged groups and waters in hydrophobic pockets.
    \item Ligand Strain Energy (Eligand•strain): This penalises poses with close internal contacts within the ligand. It is a challenging aspect due to the rigid-receptor approximation. Still, Glide includes functions to count intramolecular heavy atom contacts and evaluate penalties based on the size and location (e.g., peripheral vs. central) of contacting groups.
    \end{enumerate}
    \item Emodel
    \begin{enumerate}
    \item Emodel is a composite energy score used for the ultimate selection of the best-docked pose.
    \item It combines the GlideScore, the ligand-receptor molecular mechanics interaction energy, and the ligand's internal strain energy (for flexible docking). This composite function is found to be superior in selecting the correct pose compared to using molecular mechanics energy or GlideScore alone.
    \end{enumerate}
\end{enumerate}


\subsection*{AUTODOCK}
AutoDock is an automated computational procedure primarily designed to predict the bound conformations of small, flexible ligands to macromolecular targets of known structure. The method balances the need for a robust and physically relevant procedure with the requirement to keep computational demands at a feasible level. Its core operation combines efficient conformation searching with rapid energy evaluation.
Here's a detailed explanation of how AutoDock searches for ligand conformations within the receptor pocket:
\begin{enumerate}
\item Conformational Searching Methods
AutoDock has evolved its search methods to handle the complex problem of exploring the vast number of possible positions, orientations, and conformations of a ligand.
\begin{enumerate}
\item Simulated Annealing (SA)
    \begin{enumerate}
    \item Initial Approach: Earlier versions of AutoDock (such as version 1.0) primarily used the Metropolis technique, also known as simulated annealing, for conformational and positional searching.
    \item Mechanism: In this method, the trial molecule performs a random walk in the space around the static target protein. At each time step, small random displacements are applied to the ligand's rigid body translation, rotation, and each torsion angle. The interaction energy for this new state is evaluated, and the new state is immediately accepted if the energy is lower. If the new energy is higher, the step is accepted probabilistically, depending on a user-defined "temperature". High temperatures allow nearly all steps to be accepted, while lower temperatures enforce stricter conformational selection.
    \item Refinement: The binding conformation is successively refined as the simulation proceeds by gradually reducing the temperature. The simulation is broken into cycles at a constant temperature, and the temperature is reduced at the beginning of each new cycle. Each cycle typically begins with the trial molecule in the lowest energy conformation found in the previous cycle.
    \item Limitations: While SA performed well for ligands with roughly eight or fewer rotatable bonds, problems with more degrees of freedom rapidly became intractable. It could also sometimes fail to reproduce crystal structures in challenging cases.
    \end{enumerate}
\item Genetic Algorithms (GA) and Lamarckian Genetic Algorithm (LGA)
    To overcome SA's limitations, AutoDock version 3.0 introduced new search methods, including a traditional Genetic Algorithm (GA) and a novel, adaptive global/local search method called the Lamarckian Genetic Algorithm (LGA). GAs use principles from natural genetics and biological evolution. A ligand's arrangement (translation, orientation, conformation) is defined by state variables, which correspond to "genes" in a chromosome (the genotype). The atomic coordinates are the "phenotype". The "fitness" is the total interaction energy of the ligand with the protein. New individuals (offspring) are generated through crossover (inheriting genes from parents) and mutation (random changes to genes). Selection is based on fitness, allowing better-suited solutions to reproduce. Real-valued genes are used for ligand translation (three Cartesian coordinates), orientation (four variables defining a quaternion), and each torsion angle. Quaternions are used for orientation to avoid "gimbal lock". New random number generators are introduced to ensure results can be reproduced on any hardware platform given the same seed values.
    
\item Lamarckian Genetic Algorithm (LGA)
    This is a hybrid search method combining the GA with a local search (LS) method. In LGA, environmental adaptations (local optimization by LS) of an individual's phenotype are "reverse transcribed" into its genotype and become heritable traits. This means the results of the local search update the individual's genotype, influencing future generations. The LS method is based on Solis and Wets, which has the advantage of not requiring gradient information about the local energy landscape, thus facilitating torsional space search. It is also adaptive, adjusting its step size based on success or failure in finding lower energies. The step size can be different for each type of gene (e.g., translational vs. rotational).  LGA has shown enhanced performance relative to SA and GA alone. Across various test systems, LGA was found to be the most efficient (finding lowest energy for a given number of evaluations), most reliable (finding reproducible conformations), and most successful (reproducing crystal structures most often) among the three search methods. For instance, LGA successfully docked systems with up to 10 rotatable bonds where SA failed.
    
\end{enumerate}

\item Exploration of Degrees of Freedom
AutoDock's search algorithms explore a ligand's conformational space by manipulating its degrees of freedom:
\begin{enumerate}
\item Rigid Body Motion: The ligand can translate (move in X, Y, Z space) and rotate (change its orientation) as a rigid body. AutoDock uses a quaternion formalism for rigid body rotations to avoid "gimbal lock".
\item Torsional Flexibility: A key strength of AutoDock is its ability to handle flexible ligands with multiple torsional degrees of freedom. Each rotatable bond in the ligand is treated as a degree of freedom, allowing for changes in the ligand's conformation. The number of torsional degrees of freedom affects the total conformational entropy lost upon binding, which is accounted for in the force field.
\item Reflective Walls: The Cartesian search space is bounded by "reflective' walls". If a ligand attempts to move outside the defined grid volume, it is "bounced" back into the relevant simulation space by immediately giving it twice the opposite translational perturbation. This prevents the ligand from wasting simulation time exploring artificial surfaces at the grid edges.
\item Flexible Rings (Advanced): AutoDock can indirectly manage flexible rings by converting them to an acyclic (open) form and using special "G" atom types along with a 12-2 pseudo-Lennard-Jones potential to restore the cyclic structure during the calculation while exploring ring conformations. This potential guides the "edge atoms" next to each other, effecting ring closure. Multiple flexible rings can be docked simultaneously using different carbon atom types (G, J, Q, GA).
\item Protein Flexibility (Limited): While the protein is typically treated as a static target to enable rapid grid-based energy evaluation, AutoDock4 incorporates explicit conformational modeling of specified sidechains in the receptor. Selected sidechains can rotate around torsional degrees of freedom, similar to how the ligand's flexibility is treated. This feature is crucial for systems where receptor motion upon binding is significant.
\end{enumerate}
\item User Input Minimization and Tools
The researcher's input is minimized; they generally only need to specify the region of the protein containing the binding site and place the ligand in an arbitrary conformation and orientation nearby. The AutoDockTools (ADT) graphical user interface assists in preparing coordinate files (creating PDBQT files with polar hydrogens, partial charges, atom types, and torsional information), designing experiments (identifying active sites, setting search parameters), and analyzing results (clustering, visualization). ADT helps define rotatable bonds in the ligand and can even select specific protein sidechains to be treated as flexible.
In essence, AutoDock searches for ligand conformations by performing numerous iterations of generating trial ligand positions and conformations (using SA, GA, or LGA), rapidly evaluating their energy with pre-computed grids, and then accepting or rejecting these trials based on the energy landscape and search algorithm's rules, ultimately aiming to find the lowest energy binding modes.
\item{enumerate}
\end{enumerate}
The detailed scoring function that Autodock uses to score and rank the docking poses.
AutoDock is designed to predict the bound conformations of small, flexible ligands to macromolecular targets of known structure. Its core operation balances computational feasibility with a robust and physically relevant procedure. This involves efficient conformation searching coupled with rapid energy evaluation.
AutoDock scores ligands after conformational search primarily using a semi-empirical free energy force field. This force field is designed to estimate the free energy of binding ($\Delta G$) of small molecules to macromolecular targets.
\begin{enumerate}
\item Grid-Based Energy Evaluation (AutoGrid): Before the actual docking simulation, a separate program called AutoGrid pre-calculates three-dimensional "affinity grids". The target protein is embedded in this grid. A probe atom is sequentially placed at each grid point, and its interaction energy with the protein is calculated and stored at that point. Separate grids are generated for each type of atom present in the ligand (e.g., carbon, oxygen, nitrogen, hydrogen). Additionally, separate grids for electrostatic potential and desolvation potential are calculated. During the docking simulation, these grids act as rapid look-up tables. For each ligand atom, its interaction energy with the protein is determined by trilinear interpolation of the affinity values from the eight surrounding grid points. This makes the energy calculation very fast, as it's proportional only to the number of atoms in the ligand, not the protein. The calculation of these grids can incorporate various levels of sophistication, from constant dielectrics to finite difference methods.
\item Force Field Components:
    ◦ The force field used by AutoDock is adapted from AMBER parameters. It accounts for various types of interactions:
    \begin{enumerate}
        \item Dispersion/Repulsion: Modeled with a Lennard-Jones 12-6 potential.
        \item Hydrogen Bonding: Treated with a directional 12-10 potential. The parameters are assigned to yield a maximum well depth of -5 kcal/mol at 1.9{\AA} for hydrogen bonds with oxygen and nitrogen, and -1 kcal/mol at 2.5{\AA} for those with sulfur. A function E(t) provides directionality based on the angle (t) from ideal hydrogen-bonding geometry.
        \item Electrostatics: Evaluated using a screened Coulomb potential. Earlier versions also used a constant dielectric of 40.
        \item Desolvation: Based on the volume of atoms that surround a given atom and shield it from solvent, weighted by a solvation parameter and an exponential term with a distance-weighting factor. The desolvation model in AutoDock 4.2 is parameterized for all supported atom types, not just carbon.
    \end{enumerate}
\item Conformational Entropy: The force field also includes a term to estimate the unfavorable entropy lost upon ligand binding due to the restriction of its conformational degrees of freedom. This term is proportional to the number of rotatable bonds in the ligand (Ntor), excluding bonds that only move hydrogens.
\item Intramolecular Energies: Intramolecular interaction energies within the ligand (dispersion/repulsion, hydrogen bonding, electrostatic) are calculated at each time step. AutoDock 4.2's default model for the unbound state assumes the ligand's conformation is the same as its bound state, making the internal energy contribution zero. Earlier versions (like AutoDock 4.0) used an extended conformation for the unbound state, but this was changed due to problems with sterically-crowded molecules. Users can also specify a custom unbound energy.
\end{enumerate}
The Scoring Function Equation
AutoDock 4.2 uses a force field that evaluates binding in two steps: first, estimating the intramolecular energetics for the transition from the unbound states of the ligand and protein to their bound conformations, and second, evaluating the intermolecular energetics of combining them in their bound conformations.
The general equation for the binding free energy ($\Delta G$) is given as:
\begin{equation} \label{eq:AUTODOCKscore-1}
\begin{split}
   \ \Delta \text{G} = \\ 
                        (\text{V}_{\text{bound}}^{\text{L-L}} - \text{V}_{\text{unbound}}^{\text{L-L}}) + \\ 
                        (\text{V}_{\text{bound}}^{\text{P-P}} - \text{V}_{\text{unbound}}^{\text{P-P}}) + \\ 
                        (\text{V}_{\text{bound}}^{\text{P-L}} - \text{V}_{\text{unbound}}^{\text{P-L}}) + \\
                        \Delta \text{S}_{\text{conf}}
\end{split}
\end{equation}
Where:
\begin{itemize}
    \item $L$ refers to the ligand.
    \item $P$ refers to the protein.
    \item $V^{L-L}$ represents the intramolecular interaction energy within the ligand.
    \item $V^{P-P}$ represents the intramolecular interaction energy within the protein.
    \item $V^{P-L}$ represents the intermolecular interaction energy between the protein and the ligand.
    \item $\Delta S_{conf}$ is the conformational entropy lost upon binding.
\end{itemize}

Each of the pairwise energetic terms ($V$) in the equation above includes evaluations for dispersion/repulsion, hydrogen bonding, electrostatics, and desolvation. These are combined with weighting constants ($W$) that have been optimized empirically based on a set of experimentally determined binding constants.
The specific formulation of these pairwise terms is given as:
\begin{equation} \label{eq:AUTODOCKscore-2}
\begin{split}
    \ V =  \\ 
        \text{W}_{\text{vdw}} \sum_{i,j} \left( \frac{A_{ij}}{r_{ij}^{12}} - \frac{B_{ij}}{r_{ij}^{6}} \right) + \\
        \text{W}_{\text{hbond}} \sum_{i,j} E(t) \left( \frac{C_{ij}}{r_{ij}^{12}} - \frac{D_{ij}}{r_{ij}^{10}} \right) + \\
        \text{W}_{\text{elec}} \sum_{i,j} \frac{q_i q_j}{\epsilon(r_{ij})r_{ij}} +  \\
        \text{W}_{\text{sol}} \sum_{i,j} (S_i V_j + S_j V_i) e^{(-r_{ij}^2 / \sigma^2)}
\end{split}
\end{equation}
Where:
\begin{enumerate}
    \item $W_{vdw}$, $W_{hbond}$, $W_{elec}$, $W_{sol}$ are the weighting constants determined through linear regression analysis. For example, in one model, dispersion/repulsion energies were weighted by 0.1485, electrostatics by 0.1146, and the desolvation term by 0.1711.
    \item The first term is a Lennard-Jones 12-6 potential for dispersion/repulsion, where $A_{ij}$ and $B_{ij}$ are parameters, and $r_{ij}$ is the distance between atoms $i$ and $j$.
    \item The second term is a directional 12-10 hydrogen bond potential, where $C_{ij}$ and $D_{ij}$ are parameters, and $E(t)$ provides directionality based on the angle $t$ from ideal H-bonding geometry.
    \item The third term is a screened Coulomb potential for electrostatics, where $q_i$ and $q_j$ are partial charges, and $\epsilon(r_{ij})$ is a distance-dependent dielectric function.
    \item The final term is a desolvation potential based on the volume of atoms ($V$) that surround a given atom, weighted by a solvation parameter ($S$) and an exponential term with a distance-weighting factor $\sigma = 3.5${\AA}. The summation for this term in certain models, like Model C, is performed only over carbon atoms in the ligand.
\end{enumerate}
Additionally, the torsional free energy ($\Delta \text{G}_{\text{tor}}$) is included as a separate term, typically proportional to the number of rotatable bonds ($\text{N}_{\text{tor}}$) in the ligand, with a coefficient determined during calibration (e.g., 0.3113 kcal/mol per torsion). The full set of empirically determined coefficients for these terms were derived using linear regression against known binding constants from a set of protein-ligand complexes.
After the conformational search, AutoDock performs a cluster analysis on the lowest energy docked structures found from multiple simulation runs. These clusters are then ranked by increasing energy. The software reports the fitness (docked energy), state variables, and coordinates of the docked conformation, along with the estimated free energy of binding ($\Delta\text{G}_{\text{pred}}$)


\subsection*{AUTODOCK-VINA}



\subsection*{FLEXX}

\subsection*{SURFLEX}

\subsection*{GOLD}

\printbibliography
%\bibliography{Molecular_Docking_Methods_and_Programs}
\end{document}